\documentclass[11pt,a4paper]{article}
\usepackage[margin=2.2cm]{geometry}
\usepackage{times}
\usepackage[T1]{fontenc}
\usepackage{parskip}
\setlength{\parskip}{0.5em}
\pagestyle{empty}
\usepackage{titlesec}
\titlespacing*{\section}{0pt}{0.9em}{0.3em}
\titleformat{\section}{\bfseries\large}{}{0pt}{}

\begin{document}

\begin{center}
{\Large\bfseries What structure-prediction models have learned about GPCR activation}\\[3pt]
{\small Aditya Sengar \quad$\cdot$\quad Progress summary}
\end{center}

\section*{What I am doing}

I am testing what modern co-folding models have actually learned about G-protein-coupled receptor
activation. The analysis covers around forty receptors for which reliable active and inactive
reference structures exist, predicted by five models: Boltz-2, OpenFold-3, Protenix, Chai-1 and
AlphaFold2-multimer.

The approach is to give each receptor a different binding partner and ask what the model does with
it. The partners span a deliberate range: the receptor's own cognate G protein, a scrambled or
truncated version of that protein, a designed synthetic peptide, a drug-like ligand, or no partner
at all. Each resulting structure is scored on the two established structural hallmarks of
activation, namely the outward displacement of transmembrane
helix 6 and the NPxxY microswitch on helix 7, and classified as active or inactive against that receptor's own crystal references.
By comparing across partners we can separate a model that understands receptor activation from one
that has simply learned to associate a G-protein-shaped object with an active-looking structure.

\section*{Results}

\textbf{The activation score checks out against experiment.} Everything that follows depends on the
score being meaningful, so that was established first. The active/inactive call uses only the helix 6
measurement, which leaves the NPxxY microswitch as a completely independent test of whether the call
is right. It mostly passes: 83 per cent of structures labelled active have an NPxxY distance within
1.5~\AA\ of that receptor's own active crystal. Two separate structural measurements agree with each
other and with the experimental structures, neither having been tuned to match the other, so the
classification is capturing something real rather than reflecting one arbitrary cutoff.

\textbf{Where the two measurements disagree is a finding in itself.} In 12 per cent of active calls
the NPxxY distance is in fact closer to the inactive crystal than the active one, meaning helix 6 has
swung outward while the microswitch stayed shut. This is not scattered noise. It affects 54 per cent
of active calls for dopamine D2, 43 per cent for EDNRB and 37 per cent for HRH3, but only 9 per cent
for the $\beta_2$-adrenergic receptor and none at all for a dozen others. Since activation is
understood as a coordinated motion of both elements, a model producing one half without the other is
worth reporting on its own. In practical terms the helix 6 measurement overstates activation for a
specific group of receptors, and any claim about those receptors needs the second measurement
reported alongside it.

\textbf{Partner presence drives activation.} The baseline effect is large and consistent. With its
cognate G protein present, a receptor is predicted in the active state 94 per cent of the time;
with no partner at all, 12 per cent. When the comparison is restricted to conditions where several
models were run on the same receptors, all four give essentially the same answer: between 93 and
95 per cent active with the cognate partner, and between 5 and 22 per cent without it. The effect
is therefore robust and not specific to any one implementation.

\textbf{But identity matters less than expected, and position matters more.} The obvious follow-up
is whether the model is reading the partner's sequence or merely its shape. This was tested with
designed peptides built on an inert poly-alanine backbone. A bare poly-alanine peptide never
activates the receptor (0 of 259 predictions). The same peptide carrying three substituted residues
still fails when those residues sit at arbitrary positions (4 per cent, n\,=\,257). But place the
identical three residues at positions 15, 19 and 21 and the receptor is driven active in 75 per cent
of predictions (n\,=\,525). Amino-acid composition is held constant throughout; only position
changes. The model has learned a specific positional signature rather than a general preference for
peptide bulk or hydrophobicity.

\textbf{Different models have learned different rules.} This positional signature is not a shared
property of co-folding models. Given exactly the same position-matched peptides, Boltz-2 produces an
active receptor 85 per cent of the time (n\,=\,465) while Protenix does so in none of its runs
(n\,=\,60). More broadly, across 156 conditions run on more than one model, the rank correlation
between models ranges only from 0.50 to 0.78, with OpenFold-3 and Protenix
agreeing most closely and Boltz-2 and Chai-1 least. A finding established on one model cannot be assumed to transfer to
another. AlphaFold2-multimer is included in the corpus but is not decisive at this stage: its runs
cover only ligand-bound conditions and it is never paired with a G protein, so it cannot yet
contribute to these comparisons.

\textbf{There is a sharp length threshold.} Truncating the partner peptide gives a clean
dose-response. Fragments of 5 and 10 residues essentially never activate (0 and 2 per cent), 15
residues activates a third of the time, and 20 and 25 residues activate almost always (83 and 100
per cent). The threshold is receptor-specific: the $\beta_2$-adrenergic receptor crosses between 15
and 20 residues, dopamine D2 crosses between 10 and 15, and serotonin 5-HT2A does not cross at any
length tested. Within each individual receptor the relationship is monotonic; what differs between
receptors is where the step falls. This suggests the partner contributes through two separable
routes: a short peptide must carry the correct residues in the correct places, whereas a longer
one appears to supply enough signal through its size and fold alone.

\textbf{How far this generalises is the main open question.} The ability to distinguish a genuine
partner from a decoy is confined to a minority of receptors. Counting only those with verified
crystal references, the $\beta_2$-adrenergic receptor separates its cognate G protein from decoy
helices almost perfectly (100 versus 4 per cent). At the other extreme FSHR and the M2 muscarinic
receptor barely separate them at all (100 versus 100, and 100 versus 95 per cent), treating almost
any helical peptide as sufficient. CXCR4, EDNRB and PD2R2 sit in between, at differences of 37 to
50 percentage points. The clean, interpretable behaviour that makes $\beta_2$ so informative is
therefore not typical of the wider panel, and this bounds how far the positional result can
currently be pushed.

\section*{Limitations and next steps}

These findings come from a corpus that was assembled incrementally rather than from a single
designed experiment, and the conditions are not evenly crossed, since some models were run on many
partner types and others on very few. Several of the results rest on one receptor or a small number
of them, most notably the ligand comparisons and the cleanest positional effects.

There is also a data-provenance gap worth closing. About a fifth of the active/inactive labels
belong to receptors that carry no reference structure in the shipped threshold table, so those calls
cannot be reproduced from the data as it stands; every one of them is uniformly a single class,
which is what one would expect from a labelling shortcut rather than a measurement. Those receptors
are excluded from the figures quoted above, and the references behind them need to be recovered
before the affected runs can be used.

A substantially larger and fully crossed set of runs will therefore be needed to validate all of
this: every model against every partner condition, across a broader receptor panel, with matched
replicate counts throughout. In particular the agonist-versus-antagonist comparison needs extending
well beyond the single receptor currently available, and the positional-code and length-threshold
results need independent replication on receptors outside the small permissive-versus-discriminating
set identified so far. Until that is done these should be treated as strong, internally consistent
observations rather than settled conclusions.

\end{document}
